
%% ============================================================
%%  PHỤ LỤC A: BẢNG PHÂN CÔNG CÔNG VIỆC
%% ============================================================
\appendix
\chapter{Bảng phân công công việc}

Trong quá trình thực tập tại Trung tâm An ninh mạng, công việc theo từng tuần được mô tả như bảng sau

\begin{longtable}{|p{1.1cm}|p{6.2cm}|p{6.2cm}|}
\caption{Phân công nhiệm vụ theo tuần}
\label{tab:task_assignment}\\
\hline
\textbf{Tuần} & \textbf{Thành viên 1 (Tùng)} & \textbf{Thành viên 2 (Phúc)} \\
\hline
\endfirsthead

\hline
\textbf{Tuần} & \textbf{Thành viên 1 (Tùng)} & \textbf{Thành viên 2 (Phúc)} \\
\hline
\endhead

\hline
\multicolumn{3}{|r|}{\textit{Còn tiếp...}} \\
\endfoot

\hline
\endlastfoot

1 & Tìm hiểu tổng quan về Trung tâm, làm quen với môi trường và cơ sở hạ tầng kỹ thuật. & Tìm hiểu quy trình làm việc, các công cụ và hệ thống hiện có tại Trung tâm. \\
\hline
2 & Nghiên cứu tổng quan về hệ thống IDS, phân loại, kiến trúc và các phương pháp phát hiện tấn công. & Nghiên cứu về Federated Learning: FedAvg, các thách thức Non-IID và client drift. \\
\hline
3 & Nghiên cứu về Federated Distillation: FedMD, FedProto, SSFL-IDS, Cronus. & Nghiên cứu về Byzantine attacks: data poisoning, model poisoning, logit poisoning. \\
\hline
4 & Nghiên cứu về Blom Robust Filter và các phương pháp robust aggregation (Krum, Median, TrimmedMean, FLAME, FLTrust). & Nghiên cứu về Energy-based OOD Detection và Evidential Knowledge Distillation. \\
\hline
5 & Thiết kế kiến trúc tổng thể của DFD-IDS, triển khai cơ chế Energy-based Vote Abstention (EVA). & Triển khai Blom Spectral Filter và Robust Logit Fusion. \\
\hline
6 & Triển khai cơ chế Decentralized Server Selection và tích hợp các module. & Triển khai Evidential Knowledge Distillation và tích hợp pipeline hoàn chỉnh. \\
\hline
7 & Chuẩn bị môi trường thực nghiệm, tiền xử lý dữ liệu CICIoT23, xây dựng kịch bản thử nghiệm. & Tiến hành thực nghiệm, thu thập kết quả trên các kịch bản: benign và poisoning attacks. \\
\hline
8 & Phân tích kết quả, xây dựng bảng so sánh với các phương pháp baseline, viết báo cáo tổng kết. & Hoàn thiện báo cáo, vẽ biểu đồ, kiểm tra lỗi và chuẩn bị tài liệu trình bày. \\
\hline
\end{longtable}

\vspace{1cm}

\pagebreak
